\documentclass[11pt,a4paper]{article}
\usepackage[margin=25mm]{geometry}
\usepackage{amsmath,amssymb,amsthm,hyperref}
\setlength{\emergencystretch}{2em}
\newtheorem{lemma}{Lemma}
\newtheorem{theorem}[lemma]{Theorem}
\newcommand{\gen}[1]{\langle#1\rangle}
\title{Completing the transitive-tail alternative}
\author{Human Compression project: bounded L044 proof component}
\date{}
\begin{document}\maketitle
Throughout $G=\gen{\rho_0,\ldots,\rho_{r-1}}\leq A_n$,
$n\geq12$, is transitive with a connected string C-marking
and every maximal omission intransitive. We assume the
alternating rank bound at all strictly smaller degrees at
least twelve. All supporting primitive, anchored-imprimitive,
small-marking and attachment theorems below are the complete
previously accepted proofs and required programs, not citations
of the present source target.

\begin{lemma}[A small tail and its anchor]\label{small}
Suppose $K=G_{>h}$ is transitive on its proper nonempty
support $X$, with $q=|X|<12$ and rank $d=r-1-h$ satisfying
$d\geq(q-1)/2$. Then $K$ is the literal five-point
$D_{10}$ of source Table~2(1), or one of its two six-point
$A_5$ markings (2),(3), in the relevant string orientation.
\end{lemma}
\begin{proof}
The action of $K$ is faithful, even and has full IP, since
it fixes the complement of $X$. Its marking is connected,
being a tail of a connected string. Every maximal omission
of $K$ is intransitive on $X$: otherwise its corresponding
global omission would already connect every point of $X$,
and adjoining the omitted generator, supported on $X$,
could not change any global omission orbit.

For $j<h$, the generator $\rho_j$ centralizes $K$ and
therefore preserves its unique nontrivial orbit $X$.
All generators except $\rho_h$ consequently preserve $X$;
global transitivity forces $\rho_h$ to take some point of
$X$ outside it. The accepted extension-anchor lemma shows
that this point is fixed by all tail generators except the
first. Thus the small anchored-imprimitive bound excludes
an imprimitive $K$ at the stated rank. Degrees below five
give no connected even transitive involutory marking:
$A_3$ has no involutions and the involutions of $A_4$
generate its commuting Klein four group. The accepted small
primitive attachment exclusion therefore leaves exactly
Table~2 rows (1)--(3), in the relevant string orientation.
\end{proof}

\begin{lemma}[The critical dihedral tail]\label{critical}
Suppose a unique eligible $i$-edge joins omission orbits $A,B$,
$h>i$ is the last shared label, $G_{<i}$ fixes $B$ when $i>0$,
and $K=G_{>h}$ is the five-point $D_{10}$ on $X\subseteq B$.
If the accepted fixed-support bounds are at equality
$n-|X|=2h+1$, the configuration is impossible for $n\geq12$.
\end{lemma}
\begin{proof}
The permutation centralizer of $D_{10}$ on five points is
trivial: centralizing its five-cycle gives a power of that
cycle, and centralizing an inverting reflection makes that
power the identity. Every earlier generator commuting with
$K$ therefore fixes $X$ pointwise.

For $i=0$, use the exact fixed paths of the accepted support
bound: $A$ is the path $a_0,\ldots,a_h$, while
$b_0,\ldots,b_{h-1}$ is the fixed prefix in $B$, followed by
$b_h\in X$. Since $\rho_0$ fixes $b_h$ and commutes with
$\rho_h,\ldots,\rho_2$, it fixes $b_1,\ldots,b_h$.
Its only other possible action besides $(a_0\ b_0)$ is on
$\{a_1,\ldots,a_h\}$. For $h\geq3$, the subgroup generated
by colours $2,\ldots,h$ is transitive on that set and its
subgroup of colours $3,\ldots,h$ has exactly one fixed point,
$a_1$. Centralizing both forces $\rho_0$ to fix the whole set,
contradicting evenness. For $h=1$ there is no nontrivial
permutation of the remaining singleton, again a contradiction.
Thus only $h=2$ could remain, but then $n=2h+6=10$.
This uses only the trivial centralizer of $K$, not its perfection.

Now take $i>0$. In the transitive-left-side case its accepted
support theorem gives $h=i+1$, $|A|=2h$ and fixed set
$A\cup\{b\}$. The involution $\tau=\rho_h$ is the full
exchange $\lambda$ of $h$ preserved pairs on $A$.
In the intransitive-left-side case, the accepted support
theorem gives the same parameters and fixed set; $A$ consists
of two faithful natural $S_h$-orbits, exchanged by $\tau$.
In both cases $B=\{b\}\cup X$, and $\tau$ moves $b$ into $X$.

Here is the exact perfect factor needed when $h$ is even.
Normalize the last two marks on $X=\{1,\ldots,5\}$ so the
last is $l=(1\ 2)(3\ 4)$ and the preceding one is
$s=(2\ 3)(4\ 5)$. Since $\tau$ commutes with $l$, its
crossing endpoint is the unique $l$-fixed point $5$.
Its restriction to $B$ is $(b\ 5)u$, with
$u\in C_{S_4}(l)$ an involution. When $h$ is even,
$\lambda$ is even, so $\tau|_B$ is even and $u$ is odd.
The only odd involutions of that centralizer are $(1\ 2)$
and $(3\ 4)$. Thus the literal projection of
$R=\gen{\tau,K}$ to $B$ is one of the accepted degree-six
$A_5$ heads, and in particular is perfect. As
$R\leq\gen\lambda\times A_5$ projects onto both factors,
perfection gives $(\lambda,1)\in R$: otherwise $R$ would
be the graph of a nontrivial homomorphism $A_5\to C_2$.
This conclusion takes no IP in a quotient.

For the transitive-left-side case, the accepted paired-block
argument gives an independent length-$h-1$ block marking of
$S_h$ for $L=G_{<i}$. If $h$ is even, its kernel in the
pair-flip space lies in the even-weight space. Every nonzero
$S_h$-invariant subspace there contains the all-one vector:
a nonconstant vector and a transposition give a weight-two
vector, whose conjugates span the even-weight space; a
nonzero constant vector already is all-one. But $\lambda\in R$
and actual IP gives $L\cap R=1$. The kernel is therefore
zero, so $L\cong S_h$. Its endpoint omission fixes a point
and has order $(h-1)!$, whereas a transitive degree-$2h$
action of $S_h$ has point stabilizer order $(h-1)!/2$.
This is impossible. For odd $h\geq5$, the accepted paired-row
argument applies unchanged: the crossing generator fixes $X$,
and its required odd residual involution centralizes the
$S_{h-1}$ action on the other $h-1$ pairs, forcing it even.
The $h=3$ even-vector norm argument gives an order-six group
with a nontrivial point stabilizer, contradicting transitivity
on six points. These are precisely the previously proved
transitive-side equality arguments; their only substituted
premise is the perfect factor just proved.

For the intransitive-left-side case put $k=h$. The accepted
crossing-normalization lemma uses only the trivial centralizer
of the tail, its support size at least five, and transitivity
of $R$ on $B$. It excludes odd $k$, and for even $k$ gives
\[
 \rho_i=(u_k\ b)\prod_{j<k}(u_j\ v_j),\qquad
 \lambda=\prod_j(u_j\ v_j).
\]
For $k\geq4$, the accepted explicit fifth-power word and its
diagonal $S_k$ conjugates put $\lambda$ in
$P=G_{\leq i}$. The perfect-factor argument above puts it
in $R=G_{>i}$, contradicting full IP $P\cap R=1$.
The remaining $k=2$ gives $n=2k+6=10$, outside the range.
This finishes both interior cases.
\end{proof}

\begin{theorem}[The full numerical transitive-tail alternative]
Suppose the unique eligible $i$-edge $ab$ joins the two
$G_i$-orbits $A,B$, some label greater than $i$ acts
nontrivially on both, and $h$ is the largest such label.
For $i>0$ assume $G_{<i}$ fixes $B$ pointwise, as in the
source's chosen interval orientation. Suppose $h<r-1$ and
$K=G_{>h}$ is transitive on its support $X\subseteq B$.
Then either $r\leq\lfloor(n-1)/2\rfloor$, or an actual
first or last three-generator subgroup is one of the two
source Table~2 six-point $A_5$ markings, allowing full
string reversal at the last end.
\end{theorem}
\begin{proof}
Assume the rank bound fails. Put $q=|X|$, $d=r-1-h$ and
$F=n-q$. The complete accepted support bounds give
\[
                            F\geq2h+1.                 \tag{1}
\]
For $i=0$ this is the fixed-prefix bound. For $i>0$ it is
the transitive-side bound if $G_{<i}$ is transitive on $A$,
and the completed intransitive-side bound otherwise.
Here $B$ is nontrivial because it contains $X$.
Their exact equality conclusions will be used below.

The proper support $X$ supplies an endpoint anchor exactly
as in Lemma~\ref{small}; the same proof gives connectedness,
evenness and intransitivity of all tail omissions.
If $q<12$ and $d<(q-1)/2$, then integrality and (1) give
$2r\leq q-2+2h+2\leq n-1$, a contradiction.
For the other small case Lemma~\ref{small} leaves either
an actual $A_5$ tail, giving the stated exception, or
$q=5,d=2$. In the latter case rank failure and (1) force
$F=2h+1$; Lemma~\ref{critical} excludes it. If $q\geq12$ and $K$ is proper primitive or
imprimitive on $X$, the accepted primitive or anchored
imprimitive theorem gives $d\leq\lfloor(q-2)/2\rfloor$.
The same inequality again contradicts rank failure.

Consequently only the full natural $K=A_X$ with $q\geq12$
remains. The explicit smaller-degree induction gives
$d\leq\lfloor(q-1)/2\rfloor$. Together with (1) and the
failure of the desired bound this forces precisely
\[
 q\text{ odd},\qquad d=(q-1)/2,\qquad
 F=2h+1,\qquad n=2r.                                   \tag{2}
\]
For $i=0$, the accepted exact alternating-tail endpoint
lemma applies to the equality fixed set. It gives $h=2$
and the actual first triple Table~2(2), with no decoration
on $X$.

For $i>0$ with $G_{<i}$ transitive on $A$, its support
theorem gives $h=i+1$, $|A|=2h$ and fixed set
$A\cup\{b\}$. The complete transitive-side alternating
equality obstruction excludes this configuration.

For $i>0$ with $G_{<i}$ intransitive on $A$, its completed
support theorem makes the two faithful orbit sizes equal
to $k=i+1=h$, with fixed set $A\cup\{b\}$ and diagonal
$S_k$ action on $A$. The accepted exact alternating-tail
reduction excludes $k\geq3$ and leaves $k=2$: the first
two marks form the pure degree-five $D_{10}$ on
$A\cup\{b\}$. The next mark exchanges the two pairs in
$A$ and exchanges $b$ with one point of $X$; its remaining
action is supported on $X$. Thus this is again an actual
supported dihedral attachment, now at the first end.
The complete six-choice adapter gives the rank bound or
the claimed first $A_5$ triple. This accounts for every
case under (2) and proves the theorem.
\end{proof}

\paragraph{Source and precise boundary.}
The theorem completes the numerical alternative of
Cameron--Fernandes--Leemans--Mixer, \emph{Highest rank of
a polytope for $A_n$}, arXiv:1605.09173v1,
Proposition~5.26, physical page~34, with its unique-crossing
and chosen interval orientation explicit. Lemmas~1--2 supply the missing small dihedral-tail
case by an actual anchor and a perfect-factor argument.
A single exit from the five-point tail is \emph{not} asserted
to make the entire neighboring omission have two orbits;
other eligible edges may lie in its complement. The source's rank-deficit strategy and its results
are retained in attribution; the equality obstructions and
complete dihedral adapter are the present replacement proof.
All complete supporting manuscripts, finite programs and
retained Whiston/classification premises remain paid.
There is no new census or formal theorem, and no assumption
of IP for a projected group. The intransitive-tail argument
of Proposition~5.27, later source tables, connected-square
rigidity and whole L044 recovery remain separate; no new
original or question credit is claimed.
\end{document}
